\chapter{Conclusion}
\label{ch:conclusion}
The purpose of this project is to design and implement an AI-based wardrobe assistant that supports outfit recommendation, virtual try-on, and wardrobe management. After nearly six months of research and development, the system has finally achieved most of the core functions defined in the requirements, including clothing upload and management, structured output recommendation based on LLM, and virtual try-on process. At the same time, we have also extended and implemented a calendar module to record daily clothing and provide wardrobe analysis functions, providing statistics and suggestions for users' dressing habits and wardrobe structure. 

Through collaboration between front-end interfaces, back-end APIs, and artificial intelligence components, we have built a complete workflow from user input to result display. Some key design decisions, such as adopting a more stable cloud model and combining lightweight local solutions, are determined in the early stages, while other decisions, such as introducing structured responses to enhance front-end display, are gradually refined during the development process. This indicates that the system is gradually improving through continuous trial and error and iteration.

Although the system has met most of the core FRs, there is still room for improvement. For example, at the functional level, the current system is still mainly used by individual users, and in the future, user experience and functional integrity can be further enhanced. On the one hand, emotional interaction mechanisms such as voice interaction, personalized memory, and recommendation feedback that is closer to user habits can be introduced to make the system more intelligent and user-friendly. On the other hand, social functions can be expanded, such as modules like "Inspiration Sharing" and "Shared Wardrobe", allowing users to communicate their fashion ideas, thereby enhancing the system's sense of participation and practicality. In addition, the current recommendation results are still mainly based on a single scheme, and in the future, multiple schemes can be supported for users to choose from, while further enhancing the realism of virtual try on results.